\section{From the local estimate to pointwise convergence}
\label{sec:transfer}

We now deduce Theorem~\ref{thm:main} from Theorem~\ref{thm:local}.
This section is independent of the internal notation of the analytic
proof.  It follows the transference scheme of
\cite[Sections~2--3]{V}, including its use of signed smooth profiles.
We give all the details, both to handle arbitrary integer slopes and to
separate the two roles of the hypotheses: the local estimate gives
almost everywhere convergence of smooth averages on every invertible
probability preserving system; mixing identifies their limit.

Throughout, \(b_0=0,b_1,b_2,b_3\) are the fixed primitive slopes,
\(S\) is an invertible measure-preserving transformation, and
\(e_j\) are the integers in \eqref{eq:bezout}.  All constants in
this section may depend on the slopes.

\subsection{Smooth profiles and cancellation kernels}

\begin{definition}[Averaging profile]
\label{def:tr-profile}
An averaging profile is a real \(\varphi\in C_c^\infty(\R)\) with
\begin{equation}
 \int_\R\varphi=1,
 \qquad \int_\R t^j\varphi(t)\,dt=0\quad(1\leq j\leq3),
 \label{eq:tr-moments}
\end{equation}
and, for some \(G\geq1\),
\begin{equation}
 \|\varphi^{(n)}\|_\infty
 \leq \bigl(G(n+2)\bigr)^{G(n+2)}\qquad(n\geq0).
 \label{eq:tr-gevrey}
\end{equation}
Write \(\varphi_a(t)=a^{-1}\varphi(t/a)\) for \(a>0\).
Profiles are necessarily signed: a nonnegative smooth function of
mass one cannot have zero second moment.
\end{definition}

The derivative class in \eqref{eq:tr-gevrey} is closed under fixed
translations, dilations, differentiation, and finite linear combinations.
It contains a nonnegative bump \(\rho\) of mass one supported in
\([-1,1]\).  For completeness, normalize
\[
 \rho_0(t)=
 \begin{cases}
  \exp\bigl(-1/(1-t^2)\bigr),& |t|<1,\\
  0,& |t|\geq1.
 \end{cases}
\]
For \(|t|<1\), put \(d=1-t^2\).  On a complex disk of radius \(cd\)
about \(t\), with sufficiently small absolute \(c>0\), the real part
of \((1-z^2)^{-1}\) is at least \(c'/d\).  Cauchy's estimate gives
\[
 |\rho_0^{(n)}(t)|\leq n!(cd)^{-n}e^{-c'/d}.
\]
For each \(n\) this tends to zero at the endpoints, and its maximum
over \(0<d\leq1\) is at most \(C^{n+1}n!(n+1)^n\).  Thus the zero
extension is smooth and satisfies \eqref{eq:tr-gevrey} after enlarging
\(G\).  Lemma~\ref{lem:tr-density} will construct averaging profiles
arbitrarily close to interval indicators.

Fix a profile \(\varphi\), and put
\begin{equation}
 \psi=\varphi_2-\varphi,
 \qquad
 \Delta_k(t)=\varphi_{2^{k+1}}(t)-\varphi_{2^k}(t)
            =2^{-k}\psi(2^{-k}t).
 \label{eq:tr-difference}
\end{equation}

\begin{lemma}[Realizing a scale difference]
\label{lem:tr-kernel}
There exist a dyadic integer \(D\geq1\), a constant \(C_0\geq1\),
and a real function \(\mathcal W\in C_c^\infty(\R^2)\), depending
only on \(\varphi\) and the slopes, with the following property.
For an interval \(I\) of length \(r=D2^k\) and midpoint \(m_I\),
the kernel
\[
 w_I(x,t)=\mathcal W\left(\frac{x-m_I}{r},\frac{t}{r}\right)
\]
is admissible for Theorem~\ref{thm:local}, with interior margin
\(\delta_0=1/8\) and derivative constant \(C_0\).  Moreover,
\begin{equation}
 \frac1r\int_\R\mathcal W(z,t/r)\,dz=\Delta_k(t)
 \qquad(t\in\R).
 \label{eq:tr-kernel-integral}
\end{equation}
\end{lemma}

\begin{proof}
The moments of \(\psi\) through order three vanish: its \(j\)th
moment is \((2^j-1)\int t^j\varphi(t)\,dt\).  Choose \(B\geq1\)
with \(\operatorname{supp}\psi\subseteq[-B,B]\), and define
\[
 g(t)=\frac16\int_{-\infty}^t(t-v)^3\psi(v)\,dv.
\]
For \(t>B\), expansion of the cubic and the four zero moments show
that \(g(t)=0\); it also vanishes for \(t<-B\).  Hence
\(g\in C_c^\infty(\R)\), \(g^{(4)}=\psi\), and \(g\) satisfies
\eqref{eq:tr-gevrey} with a possibly larger constant.  Indeed, its
first four derivative bounds follow from integration over \([-B,B]\),
and its subsequent derivatives are derivatives of \(\psi\).

Take \(u(z)=8\rho(8z)\), with the bump constructed above.  Choose a
dyadic integer \(D>4B\max_j|b_j|\), and set
\[
 F(z,\sigma)=u(z)D^{-3}g(D\sigma),\qquad
 L_j=\partial_\sigma-b_j\partial_z,\qquad
 \mathcal W=\prod_{j=0}^3L_jF.
\]
The constant-coefficient operators commute.  On the support of
\(\mathcal W\),
\[
 |z+b_j\sigma|\leq\frac18+\frac{|b_j|B}{D}<\frac38.
\]
After rescaling, every \(x+b_jt\) therefore lies in \(I\), at distance
greater than \(r/8\) from its endpoints.  Put
\(G_j=\prod_{i\ne j}L_iF\).  The chain rule gives
\[
 w_I(v-b_jt,t)
 =r\frac{d}{dt}
   G_j\left(\frac{v-b_jt-m_I}{r},\frac{t}{r}\right).
\]
The function differentiated is compactly supported in \(t\), so its
integral is zero.  This proves all four line cancellations.  The product
rule and \eqref{eq:tr-gevrey} give the required derivative bounds,
with fixed \(D,b_j\) absorbed into \(C_0\).

Finally, integration in \(z\) kills every term of \(\prod_jL_jF\)
containing a \(z\)-derivative.  Since \(\int u=1\),
\[
 \int_\R\mathcal W(z,\sigma)\,dz=Dg^{(4)}(D\sigma)=D\psi(D\sigma).
\]
Division by \(r=D2^k\) yields \eqref{eq:tr-kernel-integral}.
\end{proof}

We use the block size and variation \(\mathcal V_B\) from
Section~\ref{sec:local}.  Block endpoints are always deterministic;
the input values may depend on the spatial variable.

\begin{lemma}[A real-line test]
\label{lem:tr-real}
Let \(0\leq k_-\leq k_+\), and partition
\(\{k_-,\ldots,k_+\}\) into \(w\) nonempty consecutive blocks.
Suppose that bounded Borel functions \(h_1,h_2,h_3\) and
\(h_{0,k}\) are supported in a common interval of length \(L>0\),
that \(|h_j|\leq1\) for \(1\leq j\leq3\), and that
\(\mathcal V_{B_0}(h_0)\leq w^{-1/2}\) on every block \(B_0\).
Then
\begin{equation}
 \left|\sum_{k=k_-}^{k_+}\int_{\R^2}
 h_{0,k}(x)\prod_{j=1}^3h_j(x+b_jt)\Delta_k(t)\,dx\,dt\right|
 \leq C_{\varphi,b}(L+D2^{k_+}).
 \label{eq:tr-real}
\end{equation}
The constant is independent of the scale range, partition, and inputs.
\end{lemma}

\begin{proof}
Put \(r_k=D2^k\) and \(R=r_{k_+}\).  For \(0\leq a<R\), use
the translated dyadic lattice
\[
 \mathcal D_a
 =\{[a+m2^q,a+(m+1)2^q):m,q\in\Z\}.
\]
At scale \(r_k\), use the kernels from Lemma~\ref{lem:tr-kernel}.
Only roots of length \(R\) meeting the common support interval can
contribute: the zeroth coordinate of a supported kernel belongs to
its interval.  These roots have total length at most \(L+2R\).

Restrict all inputs to each such root.  This changes no local form,
since all four coordinates on its kernel support belong to the
summation interval.  In the root tree, depth is \(s=k_+-k\).
Reversing the scale order preserves the block size and variation.
Thus slot zero satisfies the schedule hypothesis of
Theorem~\ref{thm:local}, and the three fixed slots satisfy it with
one block each.  Summing that theorem over the roots gives
\begin{equation}
 \left|\sum_{k=k_-}^{k_+}
    \sum_{\substack{I\in\mathcal D_a\\|I|=r_k}}
    r_k H_I(h_{0,k},h_1,h_2,h_3)\right|
 \leq C_{b,1/8,C_0}(L+2R).
 \label{eq:tr-root-sum}
\end{equation}
For \(r=r_k\), the following sum is \(r\)-periodic in \(a\).
Since \(r\) divides \(R\), unfolding one period yields
\begin{align*}
 \frac1R\int_0^R
  \sum_{\substack{I\in\mathcal D_a\\|I|=r}}r^{-1}w_I(x,t)\,da
 &=\frac1r\int_0^r\sum_{m\in\Z}r^{-1}
   \mathcal W\left(\frac{x-a-(m+1/2)r}{r},\frac tr\right)\,da\\
 &=\frac1r\int_\R\mathcal W(z,t/r)\,dz
 =\Delta_k(t).
\end{align*}
As \(rH_I\) is the integral against \(r^{-1}w_I\), averaging
\eqref{eq:tr-root-sum} in \(a\) gives \eqref{eq:tr-real}.
All exchanges of integrals and sums involve finitely many scales,
bounded supports, and uniformly finitely many relevant intervals.
\end{proof}

\subsection{Integer sequences and finite orbit segments}

The next step is where primitivity is useful.  We do not assume that
any pair of slopes has difference one.

\begin{lemma}[An integer-sequence test]
\label{lem:tr-integer}
Under the scale and block hypotheses of Lemma~\ref{lem:tr-real},
let \(h_1,h_2,h_3\) and \(h_{0,k}\) instead be sequences on \(\Z\),
supported in a common integer interval of cardinality \(L\geq1\).
Assume the same pointwise size and block variation bounds.  Then
\begin{equation}
 \left|\sum_{k=k_-}^{k_+}\sum_{m,n\in\Z}
 h_{0,k}(m)\prod_{j=1}^3h_j(m+b_jn)\Delta_k(n)\right|
 \leq C'_{\varphi,b}(L+D2^{k_+}).
 \label{eq:tr-integer}
\end{equation}
\end{lemma}

\begin{proof}
Write \(E_b=\sum_j|e_j|\), and fix \(0<\rho<1/4\) so small that
\begin{equation}
 \rho\bigl(2+\max_j|b_j|E_b\bigr)<1.
 \label{eq:tr-radius}
\end{equation}
Plant every sequence value on the half-open interval
\([m-\rho,m+\rho)\), and set it equal to zero elsewhere.  The
resulting Borel functions have the same size and variation bounds,
and their common real support lies in an interval of length at most
\(L-1+2\rho\leq L\).

Suppose a configuration lies in four planted intervals, with integer
centers \(p_j\), and write
\(\delta_j=x+b_jt-p_j\), so \(|\delta_j|<\rho\) away from null
boundaries.  Set
\[
 m=p_0,\qquad n=\sum_je_jp_j,\qquad
 \xi=x-m,\qquad\eta=t-n.
\]
By \eqref{eq:bezout}, \(\eta=\sum_je_j\delta_j\).  Consequently
\[
 |p_j-(m+b_jn)|
 \leq |\delta_j|+|\xi|+|b_j\eta|
 <\rho(2+|b_j|E_b)<1.
\]
The expression inside the absolute value is an integer, so
\(p_j=m+b_jn\) for all four slots.  Conversely, a configuration
lies in the intervals with these centers precisely when
\[
 (\xi,\eta)\in P_\rho
 :=\{(\xi,\eta)\in\R^2:|\xi+b_j\eta|<\rho\ (0\leq j\leq3)\},
\]
up to boundaries.  This polygon has positive finite area
\(A_\rho\): it contains a neighborhood of zero and is bounded because
two distinct coordinate forms are linearly independent.  Distinct
integer pairs \((m,n)\) give disjoint regions up to boundaries,
by uniqueness of the centers and \eqref{eq:bezout}.  On
\(P_\rho\), \(|\eta|\leq\rho E_b\).

It follows that the real-line integral at scale \(k\) is exactly
\[
 I_k=\sum_{m,n\in\Z}h_{0,k}(m)\prod_{j=1}^3h_j(m+b_jn)
       \int_{P_\rho}\Delta_k(n+\eta)\,d\xi\,d\eta.
\]
Let \(D_k\) be the discrete sum at scale \(k\) in
\eqref{eq:tr-integer}.  The derivative bound
\(\|\Delta_k'\|_\infty=2^{-2k}\|\psi'\|_\infty\) gives
\[
 |\Delta_k(n+\eta)-\Delta_k(n)|
 \leq \rho E_b\|\psi'\|_\infty2^{-2k}.
\]
If this difference is nonzero, then
\(|n|\leq B2^k+\rho E_b\).  Since \(k\geq0\), there are
\(O_{\varphi,b}(2^k)\) such integers \(n\), and at most \(L\)
contributing integers \(m\).  Every coefficient product has absolute
value at most one.  Hence
\[
 |I_k-A_\rho D_k|\leq A_\rho C_{\varphi,b}L2^{-k}.
\]
Sum over \(k\), use \(\sum_{k\geq0}2^{-k}=2\), and apply
Lemma~\ref{lem:tr-real} to the planted inputs.  Division by
\(A_\rho>0\) proves \eqref{eq:tr-integer}.
\end{proof}

For measurable functions \(|f_j|\leq1\) on the probability space,
put
\begin{equation}
 F_n(x)=\prod_{j=1}^3f_j(S^{b_jn}x),\qquad
 V_k^\varphi(x)=\sum_{n\in\Z}\varphi_{2^k}(n)F_n(x)
 \quad(k\geq0).
 \label{eq:tr-orbit-averages}
\end{equation}
All sums are finite.  If
\(\operatorname{supp}\varphi\subseteq[-B_\varphi,B_\varphi]\),
with \(B_\varphi\geq1\), then
\begin{equation}
 |V_k^\varphi(x)|\leq(2B_\varphi+1)\|\varphi\|_\infty.
 \label{eq:tr-uniform}
\end{equation}

\begin{proposition}[Finite orbit oscillation]
\label{prop:tr-oscillation}
For every deterministic list \(0\leq k_1<\cdots<k_{w+1}\),
\begin{equation}
 \sum_{i=1}^w\int_X
 \max_{k_i\leq\ell\leq k_{i+1}}
       |V_\ell^\varphi-V_{k_i}^\varphi|\,d\mu
 \leq C''_{\varphi,b}\sqrt w.
 \label{eq:tr-oscillation}
\end{equation}
This holds on every invertible probability preserving system, with a
constant independent of that system, the functions, and the endpoints.
\end{proposition}

\begin{proof}
First take sequences \(h_j\), bounded by one and supported in an
integer interval \(\mathcal K\) of cardinality \(L\), and set
\[
 \mathsf V_k(m)=\sum_n\varphi_{2^k}(n)
                   \prod_{j=1}^3h_j(m+b_jn).
\]
For \(m\) in a subset \(\mathcal E\subseteq\mathcal K\), let
\(\ell_i(m)\) be the smallest maximizing index of
\(|\mathsf V_\ell(m)-\mathsf V_{k_i}(m)|\) in the finite range
\(k_i\leq\ell\leq k_{i+1}\).  Write
\(d_i(m)=\mathsf V_{\ell_i(m)}(m)-\mathsf V_{k_i}(m)\), and take
\(\alpha_i(m)=\overline{d_i(m)}/|d_i(m)|\) when \(d_i(m)\ne0\),
and zero otherwise.  In the block \(k_i\leq k<k_{i+1}\), set
\[
 h_{0,k}(m)=
 \begin{cases}
  \alpha_i(m)/(4\sqrt w),&m\in\mathcal E, k<\ell_i(m),\\
  0,&\text{otherwise}.
 \end{cases}
\]
On each block this sequence has size at most \((4\sqrt w)^{-1}\)
and at most one jump of that size.  Its block size and variation is
therefore at most \(w^{-1/2}\).  The block endpoints remain
deterministic, although the prefix and phase depend on \(m\).

The identity \(\Delta_k=\varphi_{2^{k+1}}-\varphi_{2^k}\)
telescopes, so Lemma~\ref{lem:tr-integer} gives
\begin{equation}
 \sum_{m\in\mathcal E}\sum_{i=1}^w
  \max_{k_i\leq\ell\leq k_{i+1}}
       |\mathsf V_\ell(m)-\mathsf V_{k_i}(m)|
 \leq4C'_{\varphi,b}\sqrt w\,(L+D2^{k_{w+1}}).
 \label{eq:tr-sequence-oscillation}
\end{equation}
Indeed, the sum tested in that lemma is exactly \((4\sqrt w)^{-1}\)
times the nonnegative left side, since \(\alpha_i(m)d_i(m)=|d_i(m)|\).

We now use finite orbit segments, the elementary form of Calder\'on's
transference principle \cite{Calderon}.  Fix the endpoints, and choose
an integer
\[
 Q\geq B_\varphi\max_j|b_j|\,2^{k_{w+1}}.
\]
For each \(M\geq1\) and \(x\in X\), use the sequences
\[
 h_{j,x}(r)=
 \begin{cases}
  f_j(S^rx),&1-Q\leq r\leq M+Q,\\
  0,&\text{otherwise}.
 \end{cases}
\]
Their common support has cardinality \(M+2Q\).  For
\(1\leq m\leq M\) and all scales in the endpoint list, no term
of the average is cut off, so
\(\mathsf V_k(m)=V_k^\varphi(S^mx)\).  Apply
\eqref{eq:tr-sequence-oscillation} with
\(\mathcal E=\{1,\ldots,M\}\), integrate over \(x\), and use
invariance under \(S^m\).  This gives
\[
 M\sum_{i=1}^w\int_X
  \max_{k_i\leq\ell\leq k_{i+1}}
       |V_\ell^\varphi-V_{k_i}^\varphi|\,d\mu
 \leq4C'_{\varphi,b}\sqrt w\,(M+2Q+D2^{k_{w+1}}).
\]
The maxima are measurable finite maxima.  Divide by \(M\) and let
\(M\to\infty\), keeping the endpoints and hence \(Q\) fixed.
This proves \eqref{eq:tr-oscillation}.
\end{proof}

\begin{lemma}[Oscillation implies convergence]
\label{lem:tr-cauchy}
Let \((Z_k)_{k\geq0}\) be a uniformly bounded sequence of measurable
complex functions on a probability space.  Suppose that, for every
deterministic \(0\leq k_1<\cdots<k_{w+1}\),
\[
 \sum_{i=1}^w\int_X
  \max_{k_i\leq\ell\leq k_{i+1}}|Z_\ell-Z_{k_i}|\,d\mu
 \leq C\sqrt w.
\]
Then \(Z_k\) converges on a measurable conull set.
\end{lemma}

\begin{proof}
We include the standard finite-block argument; compare the oscillation
methods in \cite{Bourgain,JKRW}.
The measurable function
\[
 \omega(x)=\lim_{K\to\infty}\sup_{p,q\geq K}|Z_p(x)-Z_q(x)|
\]
vanishes exactly where the sequence is Cauchy.  If its positive set
has positive measure, there are \(\delta>0\) and
\(E=\{\omega>2\delta\}\) with \(\mu(E)=d>0\).  For every \(K\),
\[
 \sup_{\ell\geq K}|Z_\ell-Z_K|
 \geq\tfrac12\omega>\delta\quad\text{on }E.
\]
By monotone convergence of the finite maxima, a deterministic
\(L>K\) satisfies
\(\int\max_{K\leq\ell\leq L}|Z_\ell-Z_K|\,d\mu\geq\delta d/2\).
Starting with \(k_1=0\), choose successive endpoints this way.
The assumed estimate would give \(w\delta d/2\leq C\sqrt w\)
for every \(w\), a contradiction.  Thus \(\omega=0\) almost
everywhere, and completeness of \(\C\) proves the claim.
\end{proof}

Proposition~\ref{prop:tr-oscillation}, \eqref{eq:tr-uniform}, and
Lemma~\ref{lem:tr-cauchy} already give almost everywhere convergence
of \(V_k^\varphi\).  Mixing has not yet been used.

\subsection{Identifying the limit by mixing}

We include the norm argument to make explicit that it accommodates
arbitrary distinct nonzero slopes and both signs of time.
The set-mixing assumption on \(S\) implies, by uniform approximation
by simple functions, that for bounded measurable \(g,h\),
\begin{equation}
 \int_Xg\,(h\circ S^m)\,d\mu
 \longrightarrow\left(\int_Xg\,d\mu\right)
                 \left(\int_Xh\,d\mu\right)
 \quad(|m|\to\infty).
 \label{eq:tr-functional-mixing}
\end{equation}

\begin{lemma}[The Hilbert space van der Corput criterion]
\label{lem:tr-vdc}
Let \((u_n)_{n\geq1}\) be bounded vectors in a complex Hilbert space,
with inner product linear in the first variable.  If the limits
\[
 \gamma_h=\lim_{N\to\infty}\frac1N
            \sum_{n=1}^{N-h}\langle u_{n+h},u_n\rangle
\]
exist for each \(h\geq1\) and tend to zero as \(h\to\infty\), then
\(N^{-1}\sum_{n=1}^Nu_n\to0\) in norm.
\end{lemma}

\begin{proof}
We record the finite estimate underlying the criterion; compare
\cite[Lemma~D.1]{HostKra}.  Extend \(u_n\) by zero outside
\(\{1,\ldots,N\}\).  The identity
\[
 H\sum_{n=1}^Nu_n
 =\sum_{m=1}^{N+H-1}\sum_{j=0}^{H-1}u_{m-j}
\]
and Cauchy's inequality give
\begin{align*}
 \left\|\frac1N\sum_{n=1}^Nu_n\right\|^2
 \leq\frac{N+H-1}{N^2H^2}\bigg(&H\sum_{n=1}^N\|u_n\|^2\\
 &+2\sum_{h=1}^{H-1}(H-h)
      \operatorname{Re}\sum_{n=1}^{N-h}\langle u_{n+h},u_n\rangle
 \bigg).
\end{align*}
If \(\|u_n\|\leq M\), first let \(N\to\infty\) with \(H\)
fixed, obtaining
\[
 \limsup_N\left\|\frac1N\sum_{n=1}^Nu_n\right\|^2
 \leq\frac{M^2}{H}+\frac2H\sum_{h=1}^{H-1}|\gamma_h|.
\]
The right side tends to zero with \(H\), since the Ces\`aro means
of \((|\gamma_h|)\) tend to zero.
\end{proof}

\begin{lemma}[Norm convergence at distinct slopes]
\label{lem:tr-norm}
If \(S\) is mixing, \(a_1,\ldots,a_r\) are distinct nonzero
integers, and \(g_1,\ldots,g_r\) are bounded measurable functions,
then
\begin{equation}
 \frac1N\sum_{n=1}^N\prod_{j=1}^rg_j\circ S^{a_jn}
 \longrightarrow\prod_{j=1}^r\int_Xg_j\,d\mu
 \quad\text{in }L^2(\mu).
 \label{eq:tr-norm}
\end{equation}
Consequently, for \(F_n\) in \eqref{eq:tr-orbit-averages} and
\(p=\prod_{j=1}^3\int_Xf_j\,d\mu\), every
\(\chi\in C_c^1(\R)\) of integral one satisfies
\begin{equation}
 \sum_{n\in\Z}a^{-1}\chi(n/a)F_n\longrightarrow p
 \quad\text{in }L^2(\mu),\qquad a\to\infty.
 \label{eq:tr-weighted-norm}
\end{equation}
\end{lemma}

\begin{proof}
Write \(U^mf=f\circ S^m\).  These operators are unitary on
\(L^2(\mu)\), with inner product
\(\langle f,g\rangle=\int f\overline g\,d\mu\).
We induct on \(r\), starting with the empty product \(r=0\).
Suppose first that one of the \(g_j\) has mean zero.  Put
\[
 u_n=\prod_{j=1}^rU^{a_jn}g_j,
 \qquad q_{j,h}=(U^{a_jh}g_j)\overline{g_j}.
\]
By invariance,
\[
 \langle u_{n+h},u_n\rangle
 =\int_Xq_{1,h}\prod_{j=2}^rU^{(a_j-a_1)n}q_{j,h}\,d\mu.
\]
For fixed \(h\), induction applies to the distinct nonzero slopes
\(a_j-a_1\), \(2\leq j\leq r\).  The average of this correlation
therefore tends to
\[
 \gamma_h=\prod_{j=1}^r\int_Xq_{j,h}\,d\mu.
\]
Here removing the final \(h\) terms and dividing by \(N\) instead
of \(N-h\) does not change the limit.  By
\eqref{eq:tr-functional-mixing}, the \(j\)th factor tends to
\(|\int g_j\,d\mu|^2\) as \(h\to\infty\).  At least one limit
is zero, so \(\gamma_h\to0\).  Lemma~\ref{lem:tr-vdc} proves
\eqref{eq:tr-norm} in the centered case.  For general inputs, write
each \(g_j\) as its mean plus a mean-zero function and expand.
The constant term is the asserted limit, and every other term
contains a centered factor, so the centered case applies.  This
completes the induction.  Only fixed \(h\) is used in the induction;
no estimate uniform in \(h\) is needed.

Apply \eqref{eq:tr-norm} to \(b_j\) and \(-b_j\).  The partial sums
\[
 R_N^\pm=\sum_{n=1}^N(F_{\pm n}-p)
 \quad\text{satisfy}\quad \|R_N^\pm\|_2=o(N).
\]
Choose \(B\geq1\) containing the support of \(\chi\), and put
\(M_a=\lfloor Ba\rfloor+1\), \(w_n^\pm=a^{-1}\chi(\pm n/a)\).
Then \(w_{M_a}^\pm=0\), and summation by parts gives
\[
 \sum_{n=1}^{M_a}w_n^\pm(F_{\pm n}-p)
 =\sum_{n=1}^{M_a-1}(w_n^\pm-w_{n+1}^\pm)R_n^\pm.
\]
Its norm is at most
\(a^{-2}\|\chi'\|_\infty\sum_{n<M_a}\|R_n^\pm\|_2=o(1)\).
To see the last assertion, split at a fixed \(N_0\), and use
\(\|R_n^\pm\|_2\leq\epsilon n\) beyond \(N_0\); the tail is
\(O_B(\epsilon)\), and the finite initial part tends to zero.
The \(n=0\) term is \(O(a^{-1})\) in norm, while the Riemann sum
\(a^{-1}\sum_n\chi(n/a)\) tends to one.  Adding the two signs and
the zero term proves \eqref{eq:tr-weighted-norm}.
\end{proof}

For each averaging profile, Lemma~\ref{lem:tr-norm} gives
\(V_k^\varphi\to p\) in \(L^2\).  Its almost everywhere limit
is therefore \(p\).  More explicitly, if \(E_\varphi\) is the
measurable conull convergence set, Fatou's lemma gives
\[
 \int_{E_\varphi}|\lim_kV_k^\varphi-p|^2\,d\mu
 \leq\liminf_k\|V_k^\varphi-p\|_2^2=0.
\]

\subsection{Removing the smooth weights}

The moment restrictions need not obstruct approximation in \(L^1\).
The correction below, from \cite[Lemma~2.8]{V}, moves a small amount
of signed mass to a large scale.  Uniform control of the supports or
derivative constants is neither asserted nor needed.

\begin{lemma}[Profiles approximate interval indicators]
\label{lem:tr-density}
For every \(c\in[1,2]\) and \(\epsilon>0\), there is an averaging
profile \(\varphi\) with
\[
 \|\varphi-h_c\|_{L^1(\R)}<\epsilon,
 \qquad h_c=c^{-1}\mathbf1_{(0,c]}.
\]
\end{lemma}

\begin{proof}
Let \(\rho\) be the mass-one bump constructed above, and set
\(\rho_\delta(t)=\delta^{-1}\rho(t/\delta)\).  For sufficiently
small \(\delta>0\), the real smooth function
\(v=h_c*\rho_\delta\) has mass one and
\(\|v-h_c\|_1<\epsilon/2\).  Indeed,
\[
 \|h_c(\,\cdot-s)-h_c\|_1=2\min(|s|,c)/c,
\]
and averaging this identity against \(\rho_\delta\) proves the
claim.  Also
\(\|v^{(n)}\|_\infty\leq\delta^{-n-1}\|\rho^{(n)}\|_\infty\),
so \(v\) satisfies \eqref{eq:tr-gevrey}.  Fix \(v\), and write
\(m_j=\int t^jv(t)\,dt\) for \(1\leq j\leq3\).

Choose four distinct real numbers \(z_0,z_1,z_2,z_3\), and define
\[
 P_j(z)=\int_\R(z+u)^j\rho(u)\,du,
 \qquad \rho_{i,A}(t)=A^{-1}\rho(t/A-z_i)
 \quad(A\geq1).
\]
The \(j\)th moment of \(\rho_{i,A}\) is \(A^jP_j(z_i)\).
Since \(P_j\) is monic of degree \(j\), the fixed matrix
\(\mathsf M=(P_j(z_i))_{j,i=0}^3\) is an invertible triangular
change of the Vandermonde matrix \((z_i^j)_{j,i=0}^3\).  Solve
\[
 \mathsf M\beta=(0,-m_1/A,-m_2/A^2,-m_3/A^3)^{\mathsf T}.
\]
The inverse matrix is fixed, hence \(\sum_i|\beta_i|=O_v(A^{-1})\).
Thus
\[
 \varphi=v+\sum_{i=0}^3\beta_i\rho_{i,A}
\]
has mass one, its first three moments vanish exactly, and its added
\(L^1\) error is \(O_v(A^{-1})\).  Take \(A\) large enough that
this error is less than \(\epsilon/2\).  Finite addition and fixed
dilations preserve \eqref{eq:tr-gevrey}, so this real compactly
supported function is the required profile.
\end{proof}

\begin{proof}[Completion of the proof of Theorem~\ref{thm:main}]
Continue with \(|f_j|\leq1\), and write
\[
 A_N(x)=\frac1N\sum_{n=1}^NF_n(x),
 \qquad p=\prod_{j=1}^3\int_Xf_j\,d\mu.
\]
For every rational \(c\in[1,2]\) and integer \(h\geq1\), choose
an averaging profile \(\varphi_{c,h}\) with
\(\|h_c-\varphi_{c,h}\|_1<2^{-h}\).  The preceding results give
a measurable conull set on which \(V_k^{\varphi_{c,h}}\to p\).
Intersect these sets over the countable family, obtaining a conull
set \(E\).

For \(N_k=\lfloor c2^k\rfloor\), the choice of endpoints in
\(h_c\) gives the exact identity
\[
 2^{-k}\sum_{n\in\Z}h_c(n/2^k)F_n(x)
 =\frac{N_k}{c2^k}A_{N_k}(x).
\]
Since \(|F_n|\leq1\), for every profile \(\varphi\),
\begin{equation}
 |A_{N_k}(x)-V_k^\varphi(x)|
 \leq\frac1{c2^k}
      +2^{-k}\sum_{n\in\Z}|h_c(n/2^k)-\varphi(n/2^k)|.
 \label{eq:tr-approximation}
\end{equation}
The second term tends to \(\|h_c-\varphi\|_1\): its summand is
a bounded compactly supported Riemann-integrable function, with
possible discontinuities only at \(0,c\).  Apply
\eqref{eq:tr-approximation} with \(\varphi_{c,h}\), then let
\(h\to\infty\).  For every \(x\in E\),
\begin{equation}
 A_{\lfloor c2^k\rfloor}(x)\longrightarrow p
 \quad\text{for every rational }c\in[1,2].
 \label{eq:tr-rational-scales}
\end{equation}

Fix an integer \(s\geq1\), and use the finite rational grid
\(c_i=1+i/s\), \(0\leq i\leq s\).  If
\(2^k\leq N<2^{k+1}\), choose the largest \(i\) such that
\(n=\lfloor c_i2^k\rfloor\leq N\).  Then
\(0\leq N-n\leq2^k/s+1\).  Averages of terms bounded by one
satisfy, whenever \(1\leq n\leq N\),
\[
 |A_N-A_n|\leq\frac{2(N-n)}{N}.
\]
By \eqref{eq:tr-rational-scales}, the maximum error over the finite
grid tends to zero for \(x\in E\).  Consequently
\[
 \limsup_{N\to\infty}|A_N(x)-p|\leq2/s\qquad(x\in E).
\]
Let \(s\to\infty\) to obtain convergence along all positive integers.

Finally undo the normalization of the functions and use
\(S=T^g\), \(b_j=\lambda_j/g\), as in the initial reduction.
If the functions are specified only as essentially bounded classes,
choose bounded representatives and remove the union of all integer
translates of their exceptional sets.  This is still a measurable
null set.  Every other exceptional set used above is measurable and
every intersection is countable.  Thus the proof needs no standardness
or completeness assumption on the probability space.
\end{proof}
